\section{Sincronização}
\label{sec:sincronizacao}

\lstinputlisting[language=Java]{generated/source/Sincronizacao.java}

\inputbytecode{Sincronizacao}

\paragraph{Lendo a listagem.}
\begin{itemize}
\item Offsets 0--3: \texttt{this} é duplicado, guardado na variável
  local 1 e usado por \texttt{monitorenter}.
\item Offsets 4--11: \texttt{contador++}, as três etapas discutidas na
  seção~\ref{sec:modelo-memoria}.
\item Offsets 14--16: \texttt{monitorexit} no caminho normal.
\item Offsets 19--23: o tratador de exceção, de tipo \texttt{any} na
  tabela, libera o monitor e relança. A segunda linha da tabela cobre o
  próprio tratador.
\end{itemize}

A JVMS (§3.14) diz que o compilador garante um \texttt{monitorexit} para
cada \texttt{monitorenter}, termine o método normalmente ou com exceção,
gerando tratadores que capturam qualquer exceção. Ela diz também que um
\emph{método} \texttt{synchronized} normalmente não usa essas instruções:
é marcado com a flag \texttt{ACC\_SYNCHRONIZED}, e a própria chamada
adquire o monitor. Por isso, num método \texttt{synchronized}, o
\texttt{javap -c} não mostra \texttt{monitorenter}; use \texttt{-v} e
procure a flag.

\paragraph{A decisão que isso muda.}
\texttt{synchronized(this)} usa como lock um objeto que qualquer código
com referência a ele também pode usar para sincronizar. Um
\texttt{private final Object lock = new Object()} deixa o lock sob seu
controle. Em JDK~21, \texttt{monitorenter} era também o ponto em que uma
virtual thread ficava presa à carrier ao bloquear; desde o JDK~24
(\citetitle{doc:jep491}) não é mais (seção~\ref{sec:virtual-threads}).

\paragraph{Ponte com Quarkus.}
O guia \citetitle{doc:quarkus-virtual-threads} registra a mesma mudança:
em Java~24 ou mais novo, blocos e métodos \texttt{synchronized} não
causam mais pinning de virtual threads. Antes
de trocar \texttt{synchronized} por \texttt{ReentrantLock} num serviço
Quarkus, confira a versão do JDK em produção.
